% ============================================================================
% MUSTERLÖSUNG: Reibungskräfte -- Zusatz (Rollreibung, Reifendruck, ABS)
% Physik Klasse 7 - Kräfte
% Stand: 05.02.2026
% ============================================================================

\documentclass[11pt,a4paper]{article}

\usepackage[utf8]{inputenc}
\usepackage[T1]{fontenc}
\usepackage[ngerman]{babel}
\usepackage{geometry}
\usepackage{helvet}
\usepackage{enumitem}
\usepackage{fancyhdr}
\usepackage{xcolor}
\usepackage{tcolorbox}
\usepackage{amssymb}
\usepackage{siunitx}
\usepackage{array}
\usepackage{textcomp}
\usepackage[hidelinks]{hyperref}

\sisetup{locale=DE, per-mode=power, inter-unit-product=\ensuremath{\cdot}}
\renewcommand{\familydefault}{\sfdefault}

\geometry{left=1.4cm, right=1.4cm, top=1.3cm, bottom=1.0cm}

% Fachfarbe Physik
\definecolor{physik}{HTML}{2c5aa0}
\colorlet{fachfarbe}{physik}
\definecolor{protokollrand}{HTML}{2a7a4b}
\definecolor{merkerand}{HTML}{b35c00}
\definecolor{loesung}{HTML}{c62828}

% Kopfzeile
\pagestyle{fancy}
\fancyhf{}
\fancyhead[L]{\small\textbf{Physik Klasse 7} | Reibungskräfte -- Zusatz}
\fancyhead[R]{\small\textcolor{loesung}{\textbf{MUSTERLÖSUNG}}}
\renewcommand{\headrulewidth}{0.4pt}
\setlength{\headheight}{14pt}

% Boxen
\tcbuselibrary{skins}

\newtcolorbox{merkkasten}[1][]{
    colback=white, colframe=fachfarbe,
    coltitle=black, colbacktitle=white,
    fonttitle=\bfseries\small,
    boxrule=1pt, arc=0pt,
    left=6pt, right=6pt, top=4pt, bottom=4pt,
    #1
}

\newtcolorbox{protokollbox}[1][]{
    colback=white, colframe=protokollrand,
    coltitle=black, colbacktitle=white,
    fonttitle=\bfseries\small,
    boxrule=1pt, arc=0pt,
    left=6pt, right=6pt, top=4pt, bottom=4pt,
    #1
}

\newtcolorbox{merkebox}[1][]{
    colback=white, colframe=merkerand,
    coltitle=black, colbacktitle=white,
    fonttitle=\bfseries\small,
    boxrule=1pt, arc=0pt,
    left=6pt, right=6pt, top=4pt, bottom=4pt,
    #1
}

\newcommand{\aufgabe}[1]{\noindent\textbf{\textcolor{fachfarbe}{#1}}}
\newcommand{\lsg}[1]{\textcolor{loesung}{\textbf{#1}}}
\newcommand{\stern}[1]{\textcolor{merkerand}{\textbf{#1}}}

\setlength{\parindent}{0pt}
\setlength{\parskip}{0pt}

\begin{document}

% ============================================================================
% TITEL
% ============================================================================
{\Large\bfseries Reibungskräfte -- Zusatz \stern{$\star\star\star$} -- \textcolor{loesung}{Musterlösung}}\\[0.2em]
\hrule
\vspace{0.4em}

% ============================================================================
% B1: ROLLREIBUNG
% ============================================================================
\begin{protokollbox}[title=B5: Rollreibung -- Warum Stahlräder bei Zügen?]
\small
\textbf{Beobachtung:} Je mehr sich ein Rad verformt (``walkt''), desto \lsg{größer} ist die Rollreibung.

\vspace{0.3em}
\renewcommand{\arraystretch}{1.3}
\begin{tabular}{|p{3cm}|p{2.5cm}|p{2.5cm}|p{4cm}|}
\hline
\textbf{Material} & \textbf{Verformung} & \textbf{$F_{\text{Roll}}$} & \textbf{Einsatz} \\
\hline
Stahlrad & sehr gering & \lsg{0,5} N & \lsg{Züge, Straßenbahnen} \\
\hline
Gummi (prall) & gering & \lsg{2,0} N & \lsg{Autos, Fahrräder} \\
\hline
Gummi (platt) & groß & \lsg{5,0} N & -- \\
\hline
\end{tabular}

\vspace{0.3em}
\textbf{Erklärung:} Züge haben Stahlräder, weil \lsg{die Rollreibung minimal ist (Energiesparen)}.

Autos haben Gummireifen, weil \lsg{sie bessere Haftung/Grip brauchen (Sicherheit)}.
\end{protokollbox}

\vspace{0.3em}

% ============================================================================
% B2: REIFENDRUCK
% ============================================================================
\begin{protokollbox}[title=B6: Reifendruck und Haftung]
\small
\renewcommand{\arraystretch}{1.3}
\begin{tabular}{|p{3cm}|p{2.5cm}|p{2.5cm}|p{3.5cm}|}
\hline
\textbf{Reifendruck} & \textbf{Kontaktfläche} & \textbf{Rollreibung} & \textbf{Haftung} \\
\hline
zu niedrig & \lsg{groß} & \lsg{hoch} & \lsg{gut (aber Verschleiß!)} \\
\hline
optimal & mittel & niedrig & gut \\
\hline
zu hoch & \lsg{klein} & \lsg{sehr niedrig} & \lsg{schlecht (Rutschgefahr)} \\
\hline
\end{tabular}

\vspace{0.3em}
\textbf{Fazit:} Der optimale Reifendruck ist ein Kompromiss zwischen \lsg{Rollreibung} und \lsg{Haftung}.
\end{protokollbox}

\vspace{0.3em}

% ============================================================================
% B3: ABS
% ============================================================================
\begin{merkebox}[title=B7: ABS -- Das Antiblockiersystem \stern{$\star\star\star$}]
\small
\textbf{Problem ohne ABS:} Wenn die Räder \lsg{blockieren}, wechselt der Reifen von \lsg{Haftreibung} zu \lsg{Gleitreibung}. Da $F_{\text{Gleit}}$ \lsg{$<$} $F_{\text{Haft}}$ ist, wird der Bremsweg \lsg{länger}!

\vspace{0.3em}
\textbf{Lösung mit ABS:} ABS erkennt das Blockieren und \lsg{verringert den Bremsdruck}, sodass das Rad wieder \lsg{rollt} und die \lsg{Haftreibung} wirkt.

\vspace{0.3em}
\textbf{Ergebnis:} Mit ABS ist der Bremsweg \lsg{kürzer} als ohne ABS, besonders auf \lsg{nassen} Straßen.
\end{merkebox}

\vspace{0.3em}

% ============================================================================
% T1: ABS MESSWERTE
% ============================================================================
\aufgabe{T3: Messwerte aus der ABS-Simulation (bei 50 km/h)}

\vspace{0.2em}
{\small
\renewcommand{\arraystretch}{1.3}
\begin{tabular}{|p{3cm}|p{3.5cm}|p{3.5cm}|p{3.5cm}|}
\hline
\textbf{Straße} & \textbf{Ohne ABS} & \textbf{Mit ABS} & \textbf{Unterschied} \\
\hline
Trocken & \lsg{15 m} & \lsg{12,5 m} & \lsg{2,5 m} \\
\hline
Nass & \lsg{45 m} & \lsg{20 m} & \lsg{25 m} \\
\hline
Vereist & \lsg{110 m} & \lsg{95 m} & \lsg{15 m} \\
\hline
\end{tabular}
}

\vspace{0.2em}
\textbf{Erkenntnis:} Der Unterschied ist am größten auf \lsg{nassen} Straßen.

\vspace{0.3em}
{\small\textcolor{fachfarbe}{\textbf{Erklärung für Lehrkräfte:}} ABS hilft am meisten auf nasser Fahrbahn, weil dort der Unterschied zwischen Haft- und Gleitreibung am größten ist. Auf Eis sind beide Reibwerte sehr niedrig, daher ist der ABS-Vorteil geringer (ABS hilft dort vor allem bei der Lenkbarkeit).}

\vspace{0.4em}

% ============================================================================
% SELBSTCHECK
% ============================================================================
\begin{merkkasten}
\textbf{Erwartete Kompetenzen nach den Zusatzaufgaben:}

\vspace{0.2em}
\small
\begin{itemize}[leftmargin=*, itemsep=2pt]
    \item Erklären, warum Züge Stahlräder haben (geringe Rollreibung)
    \item Den Zusammenhang zwischen Reifendruck und Haftung beschreiben
    \item \stern{$\star\star\star$} Die Funktionsweise von ABS erklären können
\end{itemize}
\end{merkkasten}

\end{document}
